\section{Discussion}
\label{sec:discussion}

\paragraph{How much does the drop criterion matter?} Less than its role in the loop suggests, as long as it stays close to magnitude. Replacing about a third of magnitude's drop decisions, with an entropy weighting or with a loss-based Taylor or Fisher score, changed accuracy by at most a few tenths of a point in every regime, never significantly, and for entropy-weighted magnitude, tested with five seeds, the change was within the equivalence margin in three of the four regimes. The loop is tolerant: weights dropped by mistake can be regrown by gradient at the next update, and at moderate sparsity many sub-networks of the same size train equally well. What the loop does not tolerate is a criterion that ignores magnitude. Such criteria share almost none of magnitude's choices and lose accuracy, more so as the task becomes harder: the gap between the best and the worst criterion is about three points on CIFAR-10 at 90\% sparsity but seven to eight on CIFAR-100.

\paragraph{Relation to earlier comparisons.} Our results agree with Nowak et al.\ \cite{nowak2023fantastic} that magnitude is hard to beat at low density, and that criteria which make similar decisions perform alike. One result differs: in their comparison, the SNIP-like score $|w|\cdot|g|$ fell behind magnitude at low density, whereas our Taylor score, the same product, is indistinguishable from $|w|$ at 10\% density. A likely reason, which we have not tested, is that we average $|w\,g|$ over the $\Delta T=400$ steps between updates, which removes much of the mini-batch noise of the gradient; they also report that frequent updates hurt the criteria that use gradients. Our gradient-only score $|g|$, which has no magnitude term, is clearly worse, consistent with the view that magnitude is the component that carries the drop decision.

\paragraph{Why the easy regime is a poor test bed.} On CIFAR-10 at 90\% sparsity, $|w|$ is within 0.3 points of the dense model and a random drop loses less than one point. Differences between reasonable criteria are then of the order of the seed-to-seed variation, which is where our one inconclusive result lies. Claims about drop criteria made only in this regime, including the one we set out to test, should be read with this ceiling in mind.

\paragraph{Why entropy fails as a drop score.} The failure of entropy alone is not a lack of signal but a misleading one: it is worse than a random drop in every regime. The replay offers an explanation. After a ReLU, the histogram entropy of a channel is dominated by its share of zeros, so low entropy marks the channels that are silent on most inputs. Such channels are not necessarily redundant; a channel that fires for few inputs can be the one that separates them. Entropy computed without the zeros varies little across channels and no longer changes the ranking. The original intuition, that a weight reading a near-constant input acts like a bias, holds for truly constant (dead) channels. In the replay they make up, on average per layer, 1--7\% of the input channels at 90\% sparsity and about a quarter at 98\%, but they do not need entropy to be found: a weight whose input is always zero receives no gradient and shrinks under weight decay, so magnitude drops it as well.

\paragraph{What this means for designing drop criteria.} A new drop score should be evaluated against magnitude in more than one regime, with a random-drop control to set the scale, and with a statement of how much its drop set differs from that of magnitude. A score that changes little is unlikely to matter; a score that changes a lot must earn its difference in accuracy, and activation statistics that are confounded with sparsity are a poor basis for it.

\section{Limitations}
\label{sec:limitations}

\begin{itemize}
    \item \textbf{Scale.} All regimes use $32\times32$ images (CIFAR-10/100). A Tiny ImageNet regime was prepared but not run. Rankings on CIFAR need not carry over to ImageNet \cite{deng2009imagenet}, where DST is usually evaluated \cite{evci2020rigging}.
    \item \textbf{One loop.} Regrowth is always by gradient magnitude, with one schedule and (except in Sec.~\ref{sec:allocation}) global allocation. Other regrowth rules, such as random regrowth \cite{mocanu2018scalable}, or other schedules may interact differently with the drop criterion. We did not run the reference RigL configuration (random ERK mask, no warm-up) alongside our loop.
    \item \textbf{Statistics.} Most cells have three seeds, which gives wide intervals: several loss-based criteria and the CIFAR-10 random control can be neither separated from $|w|$ nor shown equivalent. The equivalence claim rests on five seeds. Dense references have one seed. The comparison of entropy with random and the robustness checks of H4 were added after the results were seen and are labelled post hoc.
    \item \textbf{Divergence.} Drop-set overlaps are measured on the trajectory of the $|w|$ run (one seed) at two checkpoints, with the weights frozen during the replay; a criterion's own trajectory could drift further. The explanation of the entropy result by post-ReLU sparsity rests on a correlation.
    \item \textbf{Run-to-run variation.} Training is not bitwise deterministic. A first run of the C0 ablation gave $-0.04$ points for $H\cdot|w|-|w|$, the re-run $+0.27$; both are below the pre-specified threshold, and all reported numbers come from the re-run.
    \item \textbf{Hardware.} Unstructured sparsity is executed with dense kernels; we report non-zero FLOPs but no speed-ups.
\end{itemize}

\section{Conclusion}
\label{sec:conclusion}

We set out to improve the drop criterion of dynamic sparse training with activation entropy and found that it did not help. Turning the question into a controlled study of eight criteria in four regimes, with hypotheses and equivalence margins fixed in advance, gave a clearer picture. Magnitude is a necessary part of the drop score: entropy alone loses 2.6 to 7.7 points, and it is worse than dropping at random, because the histogram entropy of a post-ReLU channel mostly measures its sparsity. Entropy-weighted magnitude is equivalent to magnitude in three of four regimes and inconclusive in the easiest one, and loss-based Taylor and Fisher scores are no different from magnitude either. The loop tolerates criteria that change about a third of magnitude's decisions, but not criteria that ignore magnitude, and the cost of ignoring it grows with task difficulty. Future work should test these findings at ImageNet scale, with other regrowth rules, and with criteria that combine magnitude with an explicit estimate of the loss increase.

\section*{Code and data availability}

The code of the training loop and of every criterion, the experiment plan, the design document with the hypotheses and its dated changelog, the result files of all runs and the scripts that produce every table and figure of this paper are available at \TODO{repository URL}.
